\cotefile{79}
\scanpage{145}
% archive sheet number: 144
% L1–4
ce qui, par la définition de l'opération de $Z'_n$ sur l'une des structures $n$ polygonales sur un ensemble $E$, \deleted{\unsure{à deux élts}}, implique qu'on a un isomorphisme canonique
% L5
\[
p^{-1}(\dot\lambda)\wedge\mu(i)\wedge\mu(\dot\lambda i)\simeq
\underset{\substack{\text{torseur unité}\\\text{sur }\pm1}}{\mathbb{1}}.
\tag*{(43)}
\]
% L6–10
où
\[
 p:\underset{\substack{\Vert\\\Lambda_n^*}}{Z_n}\longrightarrow
 \underset{\substack{\Vert\\\Lambda_n^*/\pm1}}{Z'_n}
\]
\deleted{\unsure{est}} est l'hom. canonique, \deleted{\unsure{surj.}} de degré 2 (\uline{NB} on a \deleted{supposé} $n\ne2$ donc $\deleted{\unsure{+}}1\ne-1$ dans $\Lambda_n$). On a en fait que pour
% L11
\[
\lambda\in p^{-1}(\dot\lambda),\quad \omega\in\mu(i),\quad \omega'\in\mu(i'=\dot\lambda.i)
\]
% L12
on a
\[
\lambda\wedge\omega\wedge\omega'=1\quad\Longleftrightarrow\quad
\omega^{\deleted{\unsure{\lambda}}\,\prime}_E=\omega_E^\lambda
\tag*{(44)}
\]
% L13–15
où $\omega_E,\omega'_E$ sont \deleted{\unsure{interprétées comme} des} permutations circulaires de l'ensemble-étoile $E$ qui correspondent aux orientations $\omega,\omega'$.

% L16–20
A priori l'isomorphisme (43) dépend du choix du sommet $s$. Mais par transport de structure, il est le $=$ pour deux sommets $s,s'$ tels que l'on ait
\[
s'=g.s,\quad\text{avec }g\in G_0^{\unsure{*|+}}
\]
% L21–25
à condition de supposer que $g$, opérant sur $\mu(i)$ et $\mu(\unsure{\lambda'|\dot\lambda i})\unsure{\vphantom{i}^{\epsilon}}$, opère avec le $=$ signe. Si on admet que c'est tjrs le cas (cf plus bas), il suffit de prendre $g\in G_0^+$, \deleted{\unsure{et l'obstruction}}\inserted{\unsure{moy.nt|moy.} une certaine orientation} commune
\[
 G_0\xrightarrow{\mathrm{sg}}\pm1\,,\tag*{(45)}
\]
